% ==============================================================================
% CONFIGURACIÓN DE PAQUETES - LATEX
% SIGC-CUSCO | Lab 04: Modelado del Sistema: UML y Prototipos
% ==============================================================================

\usepackage[utf8]{inputenc}
\usepackage[T1]{fontenc}
\usepackage[spanish, es-tabla, es-nodecimaldot]{babel}
\spanishdeactivate{"}
\usepackage{csquotes}

\usepackage[margin=2.3cm, headheight=18pt]{geometry}

\usepackage{mathptmx}
\usepackage{setspace}
\setstretch{1.15}

\usepackage{graphicx}
\graphicspath{{./}{imagenes/}{./imagenes/}}
\usepackage{float}
\usepackage{booktabs}
\usepackage{tabularx}
\usepackage{array}
\usepackage{caption}

\usepackage[table]{xcolor}

\usepackage{enumitem}

\usepackage{fancyhdr}
\pagestyle{fancy}
\fancyhf{}
\fancyhead[L]{\small\itshape UNSAAC -- Ingeniería de Software I}
\fancyhead[R]{\small\itshape Lab 04: UML y Prototipos -- SIGC-CUSCO}
\fancyfoot[C]{\thepage}
\renewcommand{\headrulewidth}{0.4pt}
\renewcommand{\footrulewidth}{0pt}

\usepackage{xurl}
\usepackage{hyperref}
\hypersetup{
    colorlinks=true,
    breaklinks=true,
    linkcolor=blue!75!black,
    urlcolor=blue!75!black,
    pdftitle={Lab 04 - Modelado UML y Prototipos - SIGC-CUSCO},
    pdfauthor={Choquenaira, Yaranga, Ccama}
}

\usepackage{microtype}
\emergencystretch=2.5em

\setlength{\parindent}{0.6cm}
\setlength{\parskip}{0.25em}

\newcolumntype{Y}{>{\raggedright\arraybackslash}X}
\newcolumntype{P}[1]{>{\raggedright\arraybackslash}p{#1}}
